
\documentclass[10pt]{article}
\usepackage[utf8]{inputenc}
\usepackage[T1]{fontenc}
\usepackage{amsmath, amssymb, xcolor, geometry, enumitem, multicol, tcolorbox}

\definecolor{mainblue}{RGB}{0, 102, 204}
\geometry{a4paper, margin=1.2cm, top=1cm, bottom=3cm, footskip=1.2cm}

\begin{document}
\enlargethispage{2.5cm}

% Modern Header
\begin{tcolorbox}[colback=mainblue!5, colframe=mainblue, sharp corners, boxrule=0.5pt]
    \small \textbf{Unit:} undefined \hfill \textbf{Name:} \rule{3cm}{0.4pt} \\
    \small \textbf{Topic:} Variables and Expressions / Order of Operations \hfill \textbf{Date:} \rule{2cm}{0.4pt} \\
    \centering \textbf{\large Homework 3}
\end{tcolorbox}

\vspace{0.2cm}

% MCQ Section
\begin{multicols}{2}
\begin{enumerate}[label=\textbf{Q\arabic*.}, leftmargin=*, itemsep=6pt, parsep=0pt]

  \item \small What is the variable in the expression $2x + 3$?
  \begin{enumerate}[label=(\Alph*), leftmargin=0.6cm, nosep, itemsep=2pt]
    \item 2
    \item 3
    \item x
    \item 2x
    \item None of the above
  \end{enumerate}

  \item \small Identify the coefficient in the term $4y$?
  \begin{enumerate}[label=(\Alph*), leftmargin=0.6cm, nosep, itemsep=2pt]
    \item y
    \item 4
    \item 4y
    \item None of the above
    \item 0
  \end{enumerate}

  \item \small What is the constant term in the expression $x + 2$?
  \begin{enumerate}[label=(\Alph*), leftmargin=0.6cm, nosep, itemsep=2pt]
    \item x
    \item 2
    \item x + 2
    \item None of the above
    \item 0
  \end{enumerate}

  \item \small Which of the following is an example of a simple expression? $3x$
  \begin{enumerate}[label=(\Alph*), leftmargin=0.6cm, nosep, itemsep=2pt]
    \item Equation
    \item Inequality
    \item Expression
    \item None of the above
    \item Formula
  \end{enumerate}

  \item \small What operation is denoted by the symbol $+$?
  \begin{enumerate}[label=(\Alph*), leftmargin=0.6cm, nosep, itemsep=2pt]
    \item Subtraction
    \item Addition
    \item Multiplication
    \item Division
    \item None of the above
  \end{enumerate}

  \item \small Identify the operator in $7 - 2$?
  \begin{enumerate}[label=(\Alph*), leftmargin=0.6cm, nosep, itemsep=2pt]
    \item $+$
    \item $-$
    \item $	imes$
    \item $div$
    \item None of the above
  \end{enumerate}

  \item \small What does the term $3 cdot x$ represent?
  \begin{enumerate}[label=(\Alph*), leftmargin=0.6cm, nosep, itemsep=2pt]
    \item 3 divided by x
    \item 3 subtracted from x
    \item 3 multiplied by x
    \item 3 added to x
    \item None of the above
  \end{enumerate}

  \item \small Which part of the expression $2x + 5$ is the constant term?
  \begin{enumerate}[label=(\Alph*), leftmargin=0.6cm, nosep, itemsep=2pt]
    \item $2x$
    \item $5$
    \item $2x + 5$
    \item None of the above
    \item $2$
  \end{enumerate}

\end{enumerate}
\end{multicols}

\vspace{-0.2cm}
\hrule
\vspace{0.2cm}
\textbf{\small Open Ended Questions (FRQ)}
\begin{enumerate}[label=\textbf{Q\arabic*.}, start=9, itemsep=5pt, topsep=0pt]

  \item \small Draw a diagram to represent the expression $x + 2$. Label the variable and the constant. \vspace{2cm}

  \item \small Explain in your own words what a variable is in an algebraic expression. Provide an example. \vspace{2cm}

\end{enumerate}

\vfill

% Chef's Tips Footer
\begin{tcolorbox}[colback=blue!5, colframe=mainblue!50, title=\small \textbf{Chef's Tips}, fonttitle=\bfseries, bottom=1mm]
\begin{itemize}[noitemsep, topsep=2pt, leftmargin=0.5cm]
  \item \scriptsize Ensure students understand that variables can represent any value, not just numbers.\item \scriptsize Emphasize the difference between coefficients and constants in expressions.\item \scriptsize Use visual aids to help students grasp the concept of expressions and variables.
\end{itemize}
\end{tcolorbox}

\end{document}
  